\section{A labelled valuation-incidence complex}\label{sec:valuation-incidence-complex}

This section develops a higher-dimensional object directly from the prime
factorisation of an $abc$ triple.  It is deliberately restricted to
\emph{primitive nonunit triples}
\[
 a,b,c\in\N_{>0},\qquad a>1,\quad b>1,\quad a+b=c,
 \quad \gcd(a,b)=1.
\]
The restriction is part of every statement below.  Unit-arm triples require a
separate extension.  The construction is finite and exact.  The next two
sections refute its first one-face quantitative selectors and preserve the
remaining weighted, multi-face, three-arm, and homological design problem.

Put $x_A=a$, $x_B=b$, and $x_C=c$.  For $i\in\{A,B,C\}$ let
\[
 S_i=\{p:p\text{ prime},\ p\mid x_i\},\qquad e_i(p)=v_p(x_i).
\]
Primitivity and $a+b=c$ make the three coordinates pairwise coprime.  Thus
their prime supports are disjoint, although we retain the arm label because it
also records a different local congruence.

\begin{definition}[Valuation-incidence complex]
A vertex is a labelled pair $(i,p)$ with $p\in S_i$, carrying the datum
$(i,p,e_i(p))$.  A face is a triple
\[
 F=(F_A,F_B,F_C),\qquad F_i\subseteq S_i.
\]
We write $F\preceq G$ if $F_i\subseteq G_i$ on every arm.  Union and
intersection are taken armwise.  The empty face is
$\varnothing_P=(\varnothing,\varnothing,\varnothing)$ and the top face is
$\top_P=(S_A,S_B,S_C)$.
\end{definition}

The face relation is the Boolean face lattice on
$S_A\sqcup S_B\sqcup S_C$.  In particular, it is reflexive, transitive and
antisymmetric; the empty and top faces are respectively least and greatest;
and armwise union is the least common upper bound.  These assertions follow
coordinatewise from the corresponding elementary facts for finite subsets.
If
\[
 \nu(F)=|F_A|+|F_B|+|F_C|,
 \qquad \dim_{\N}(F)=\max\{\nu(F)-1,0\},
\]
then a nonempty top face with $N$ vertices has dimension $N-1$.  There is no
a priori bound on $N$.

For a face $F$ and an arm $i$, define
\[
 \begin{aligned}
 R_i(F)&=\prod_{p\in F_i}p,&
 D_i(F)&=\prod_{p\in F_i}p^{e_i(p)-1},\\
 M_i(F)&=\prod_{p\in F_i}p^{e_i(p)},&
 d_i(F)&=\sum_{p\in F_i}(e_i(p)-1).
 \end{aligned}
\]
Thus $R_i$ records selected radical mass, $D_i$ selected multiplicity beyond
the first layer, $M_i$ the corresponding prime-power modulus, and $d_i$ its
additive defect degree.  The exact six-coordinate realisation is
\[
 T(F)=((R_A(F),D_A(F)),(R_B(F),D_B(F)),(R_C(F),D_C(F))).
\]
Applying logarithms produces an additive display in $\R^6$, but logarithms
play no role in face membership or in the filtration below.

\begin{proposition}[Exact face factorisation]\label{prop:vic-face-factorisation}
For every face and every arm,
\[
 R_i(F)D_i(F)=M_i(F),\qquad M_i(F)\mid x_i.
\]
At the top face,
\[
 M_i(\top_P)=x_i,\qquad R_i(\top_P)=\rad(x_i),
 \qquad \rad(x_i)D_i(\top_P)=x_i.
\]
If $F$ and $G$ are armwise disjoint, each of $R_i,D_i,M_i$ is multiplicative
under $F\cup G$, while $d_i(F\cup G)=d_i(F)+d_i(G)$.
\end{proposition}

\begin{proof}
For every selected prime, positivity of $e_i(p)$ gives
$p\,p^{e_i(p)-1}=p^{e_i(p)}$.  Multiplication over $F_i$ proves the first
identity.  The selected product is a subproduct of the canonical
factorisation of $x_i$, so it divides $x_i$; on the top face it is the full
factorisation.  Retaining one copy of each support prime gives the usual
squarefree radical.  The disjoint-union formulas are the standard product and
sum formulas over disjoint finite sets.
\end{proof}

For a budget $u=(u_A,u_B,u_C)\in\N^3$, define the filtered subcomplex
\[
 \mathcal K_u(P)=\{F:d_i(F)\le u_i\text{ for }i=A,B,C\}.
\]

\begin{proposition}[Defect filtration]\label{prop:vic-filtration}
If $G\preceq F$ and $F\in\mathcal K_u(P)$, then
$G\in\mathcal K_u(P)$.  The empty face lies in every filtration level.  If
$F,G$ are armwise disjoint, $F\in\mathcal K_u(P)$ and
$G\in\mathcal K_v(P)$, then $F\cup G\in\mathcal K_{u+v}(P)$.  Moreover,
\[
 F\in\mathcal K_{(0,0,0)}(P)
 \quad\Longleftrightarrow\quad
 e_i(p)=1\text{ for every selected }(i,p).
\]
\end{proposition}

\begin{proof}
Every summand $e_i(p)-1$ is nonnegative.  Restricting to a subset can only
decrease its sum, and the empty sums vanish.  The graded-union assertion uses
the additive formula in Proposition~\ref{prop:vic-face-factorisation}.  At
budget zero, a finite sum of nonnegative integers is zero precisely when each
summand is zero; positivity of $e_i(p)$ then gives the final equivalence.
\end{proof}

The arm colours also retain the additive equation at selected prime-power
resolution.

\begin{proposition}[Local incidence and CRT rigidity]
Every face satisfies
\[
 c\equiv b\pmod{M_A(F)},\qquad
 c\equiv a\pmod{M_B(F)},\qquad
 a+b\equiv0\pmod{M_C(F)}.
\]
The three moduli are pairwise coprime.  Call $F$ \emph{$AB$-reconstructing}
when $c<M_A(F)M_B(F)$.  Every top face is $AB$-reconstructing.  More
generally, if $F$ is $AB$-reconstructing and
\[
 0\le x<M_A(F)M_B(F),\quad
 x\equiv b\pmod{M_A(F)},\quad x\equiv a\pmod{M_B(F)},
\]
then $x=c$.
\end{proposition}

\begin{proof}
Proposition~\ref{prop:vic-face-factorisation} gives
$M_A(F)\mid a$, $M_B(F)\mid b$, and $M_C(F)\mid c$; subtracting within
$a+b=c$ proves the three congruences.  Since each selected modulus divides
one member of a pairwise coprime triple, the moduli are pairwise coprime.
At the top face, $M_A=a$ and $M_B=b$.  For nonunits one has
$a+b\le ab$, with equality only at $a=b=2$; primitivity excludes that case,
so $a+b<ab$.
For a general reconstructing face, the displayed residues show that $x$ and
$c$ are congruent modulo both coprime moduli, hence modulo their product.
Both lie in the same half-open interval of that length, so they are equal.
\end{proof}

The top point recovers the ordinary conductor exactly:
\[
 R_A(\top_P)R_B(\top_P)R_C(\top_P)=\rad(abc).
\]
It also gives a precise half-space form of a rational-power $abc$ bound.

\begin{theorem}[Exact top-face half-space]\label{thm:vic-half-space}
Write $R_i=R_i(\top_P)$ and $D_C=D_C(\top_P)$.  For all $m,n\in\N$,
\[
 c^m\le\rad(abc)^{m+n}
 \quad\Longleftrightarrow\quad
 D_C^m\le (R_AR_B)^{m+n}R_C^n.
\]
Furthermore, for any face $F$, the stronger selected-face inequality
\[
 D_C^m\le
 \bigl(R_A(F)R_B(F)\bigr)^{m+n}R_C^n
\]
implies $c^m\le\rad(abc)^{m+n}$.
\end{theorem}

\begin{proof}
The top-face identities give $c=R_CD_C$ and
$\rad(abc)=R_AR_BR_C$.  Substitution turns the left inequality into
\[
 R_C^mD_C^m\le
 R_C^m\bigl((R_AR_B)^{m+n}R_C^n\bigr).
\]
Since $R_C>0$, cancellation proves the equivalence.  For the second claim,
$F_i\subseteq S_i$ gives $R_A(F)\le R_A$ and $R_B(F)\le R_B$; enlarge the
right side of the selected-face inequality and apply the reverse implication
of the equivalence.
\end{proof}

Theorem~\ref{thm:vic-half-space} is an exact change of coordinates.  It does
not prove that a useful selected face exists.  The following fully realised
example also prevents two tempting degeneracy assumptions.

\begin{proposition}[A five-dimensional complete-premise boundary]
The primitive nonunit triple
\[
 12+833=845,\qquad
 12=2^2\cdot3,\quad 833=7^2\cdot17,\quad 845=5\cdot13^2
\]
has a six-vertex top face, hence natural dimension five, and
\[
 T(\top_P)=((6,2),(119,7),(65,13)),\qquad
 (d_A,d_B,d_C)=(1,1,1).
\]
Every arm has nontrivial defect, and on every arm the two valuation exponents
are different.  The face selecting $3$, $17$, and $5$ lies at zero budget.
\end{proposition}

\begin{proof}
The displayed factorizations list the six vertices and give all coordinates
by their definitions.  They also show directly that one exponent on each arm
is $2$ and the other is $1$.  Thus the defects are $2,7,13$, all nontrivial,
and every defect degree is one.  Selecting the three exponent-one vertices
gives defect degree zero on every arm.
\end{proof}

This example refutes exactly the pointwise universal statements that some arm
must have defect one and that the valuations must be constant on each arm.  It
does not refute eventual or density variants and does not refute any estimate
that permits all three arms to be defective.

The first falsifiable gate proposed for this higher-dimensional route used an
absolute budget.  For fixed $m,n\ge1$, it asked whether a number
$t=t(m,n)$ exists such that all but finitely many primitive nonunit triples
possess a face $F$ satisfying
\[
 d_A(F),d_B(F)\le t,\qquad
 c<M_A(F)M_B(F),\qquad
 D_C^m\le\bigl(R_A(F)R_B(F)\bigr)^{m+n}R_C^n.
\tag{VIC-ABS-1}
\]
Theorem~\ref{thm:vic-half-space} proves that VIC-ABS-1 would imply the
corresponding rational-exponent height bound on this domain.  The next section
constructs an infinite complete-premise family that refutes this absolute
budget formulation.  That obstruction does not refute the complex or its
weighted, multi-face, three-arm, and homological replacements.

The companion Lean module
\texttt{ABCValuationIncidenceComplex20260903} constructs the actual faces with
\texttt{Nat.primeFactors} and the labels with \texttt{Nat.factorization}.  It
kernel-checks the finite-poset, factorisation, filtration, congruence, CRT,
half-space, selected-face, and witness statements above.  No selector is
inserted as an axiom or an inhabited interface.
